\paragraph{Speedup analysis.}
A large end-to-end speedup number is only credible if two things are shown
explicitly: (i) the naive baseline actually \emph{carries} that deficit, and
the deficit can be \emph{localized} to a physical cause; and (ii) the reported
gain can be \emph{decomposed} into a small number of design decisions, each of
which is separately sized. This appendix supplies both by sweeping each
design decision one at a time and reporting the objective's sensitivity
(\cref{app:deficit:sweeps}). This section is a companion to the
pre-RTL feasibility gates of \cref{app:details}: the gates certify that a design
is \emph{buildable}; this section certifies where its speedup \emph{comes from}.

\paragraph{Reading conventions (read before any number below).}
Throughout, every quantity written \textbf{MFU} in this appendix is the
\textbf{compute-floor tightness} ratio
$\rho_{\mathrm{comp}}=t_{\mathrm{floor}}/T_{\mathrm{LP}}$ of \cref{app:mfu}
(the mapper-derived compute floor over the LP/analytic latency within one
cost-model convention). It is \emph{not} model-FLOP utilization in the sense of
\cref{sec:reported_quantities}, because its numerator is a Timeloop tile
minimum carrying the reference geometry's padding rather than a useful-work
count, and it is \emph{never} a measured-silicon utilization; the searched
datapaths are hypothetical, not fabricated. \textbf{Same-silicon} comparisons
are equal-area and equal-TDP \emph{by construction} (both sides share the PE
grid, so the area/power model of \cref{app:gates} returns identical budgets).
The $97\%$ pre-fusion memory-stall figures and the naive MFU values are
\emph{naive-baseline attributions}, not properties of the co-designed schedule.
Every quantity is convention-tagged, and we never form a ratio across two
conventions. The fixed-silicon algorithmic gain on EfficientNet-B7 is
$2.21\times$ under this paper's naive convention.

% ---------------------------------------------------------------------------
\iffalse % Appendix F.1 (two-factor speedup decomposition) retained but disabled
\subsection{Two-Factor Speedup Decomposition}
\label{app:deficit:twofactor}

The headline speedup factors \emph{exactly} into two orthogonal terms:
\begin{equation}
\underbrace{\frac{T_{\mathrm{naive}}(\text{TPU-v3})}{T_{\mathrm{joint}}(\text{chip})}}_{\text{headline}}
\;=\;
\underbrace{\frac{T_{\mathrm{naive}}(\text{TPU-v3})}{T_{\mathrm{joint}}(\text{TPU-v3})}}_{f_1:\ \text{scheduling recovery}}
\;\times\;
\underbrace{\frac{T_{\mathrm{joint}}(\text{TPU-v3})}{T_{\mathrm{joint}}(\text{chip})}}_{f_2:\ \text{MAC provisioning}} .
\label{eq:twofactor}
\end{equation}
The first factor $f_1$ is a \emph{fixed-silicon}, same-area / same-TDP
algorithmic gain: it re-schedules the same $65{,}536$-MAC TPU-v3 and buys
$\mathbf{1.48}$--$\mathbf{1.78\times}$. The second factor $f_2$ holds the joint
schedule fixed and provisions the PE grid up to the co-designed chip's MAC
count, buying $\mathbf{20.0}$--$\mathbf{26.9\times}$. Their product reproduces
the per-workload Figure-2 headline (and the blueprint ``achieved speedup vs.\
naive TPU-v3'' rows) to within $1.4\%$
(\cref{tab:deficit_twofactor}); the residual is the analytic
chip-endpoint solve's utilization roll-off, cross-checked against the
blueprints' Timeloop-in-loop objective.

\begin{table}[ht!]
\centering
\caption{Two-factor decomposition of the co-design speedup
(\cref{eq:twofactor}). $f_1$ is the fixed-silicon scheduling gain on TPU-v3
(same area/TDP), $f_2$ the MAC-provisioning gain to the co-designed chip.
Reconstructed $=f_1\times f_2$; headline is the blueprint value. Convention:
config-0 naive; bf16 cycle basis; HiGHS-Optimal endpoints.}
\label{tab:deficit_twofactor}
\small
\begin{tabular}{@{}lccccc@{}}
\toprule
\textbf{Workload} & $f_1$ (sched., TPU-v3) & $f_2$ (MAC prov.)
& Reconstructed & Headline & Recon./Headline \\
\midrule
Llama-1 13B  & 1.66 & 20.31 & 33.75 & 34.22 & 0.99 \\
Llama-1 65B  & 1.73 & 19.84 & 34.33 & 34.33 & 1.000$^{\ddagger}$ \\
Llama-3 70B  & 1.48 & 26.85 & 39.84 & 40.16 & 0.992 \\
DeepSeek-V3  & 1.78 & 19.97 & 35.63 & 35.68 & 0.999 \\
\midrule
\textbf{Range} & \textbf{1.48--1.78} & \textbf{20.0--26.9}
& \multicolumn{3}{c}{\textbf{34.2--40.2} (prefill-oriented designs)} \\
\bottomrule
\end{tabular}

\vspace{2pt}
{\footnotesize $^{\ddagger}$Llama-1 65B's analytic chip-endpoint solve was
degenerate under HiGHS time limits, so its $f_2$ uses the authoritative
blueprint objective (exact by construction); the other three use the
\emph{independent} analytic chip solve and land within $1.4\%$.}
\end{table}

\paragraph{Is the provisioning factor physically usable?}
A reader is entitled to ask whether $f_2$ is real parallelism or an artifact of
a cost model in which compute time falls as $1/\mathrm{MACs}$ without bound. It
is a fair question: an array can only speed an operation up as far as that
operation supplies independent tiles to fill it, and a per-token decode GEMV
supplies few. We measured the answer directly rather than argue it. For every
event of DeepSeek-V3 we computed the \emph{saturating} MAC count
$M_{\mathrm{sat}} = (\text{independent output tiles}) \times (\text{lanes per
tile})$ under the same tile geometry the mapper uses, where independent tiles
are row tiles times column tiles (k-tiles accumulate into one output and are
not independent) times the number of GEMMs carrying work (expected active
experts for routed FFNs). Beyond $M_{\mathrm{sat}}$ a chip has no more
concurrent tiles to run for that event, and its compute time cannot keep
falling. \Cref{tab:msat} gives the distribution in multiples of the TPU-v3
array.

\begin{table}[ht!]
\centering
\small
\caption{Saturating MAC count per event, DeepSeek-V3, as a multiple of
TPU-v3's $65{,}536$ MACs. Prefill events saturate three to four orders of
magnitude above the largest chip in this paper ($24.4\times$), so the
provisioning factor $f_2$ is physically usable parallelism at prefill. Decode
events saturate at $9$--$14\times$ at the low end, which is where a bound on
usable parallelism would bind, and only for chips beyond it.}
\label{tab:msat}
\begin{tabular}{@{}llrrr@{}}
\toprule
Regime & Phase & min & p10 & median \\
\midrule
$B{=}8$, ctx $4096$   & prefill & $3{,}584\times$  & $3{,}584\times$  & $16{,}384\times$ \\
                      & decode  & $9\times$        & $14\times$       & $16\times$ \\
$B{=}32$, ctx $4096$  & prefill & $14{,}336\times$ & $14{,}336\times$ & $65{,}536\times$ \\
                      & decode  & $14\times$       & $14\times$       & $32\times$ \\
$B{=}32$, ctx $512$   & prefill & $1{,}792\times$  & $1{,}792\times$  & $8{,}192\times$ \\
                      & decode  & $14\times$       & $14\times$       & $32\times$ \\
\bottomrule
\end{tabular}
\end{table}

At the original prefill operating point, the table's tile geometry supplies
many independent matrix tiles, so useful parallelism is less restrictive than
it is for small decode matrices. The bound becomes relevant when installed
arrays exceed the independent work available to an event. These counts describe
the mapper's geometry, not placed-array timing, and do not account for growing
non-matrix costs. The native full-inference model instead charges padded matrix
service and explicit vector, local, and external service for each implementation
(Appendix~\ref{app:rows_service}); its results must be read under that contract.

\paragraph{Use only the two-factor form; do not multiply the full decision ladder.}
It is tempting to read the software-ablation ladder (tile, placement, remat,
datapath, datatype) as a chain of multiplicative factors and multiply them all.
That \emph{overshoots} the headline by $1.8$--$2.1\times$ (a mechanical product
of $70.7$--$73.2\times$ versus the $34$--$40\times$ headline), for two
convention reasons, neither of them physics:
\begin{enumerate}
\item \textbf{The INT8 throughput label double-counts.} The blueprint headline
$T_{\mathrm{joint}}$ is a \emph{bf16 cycle count} at the chip's MAC count; INT8
(W8A8) is carried as a throughput / power label (roughly half the area and power
per MAC), \emph{not} as a separate cycle factor. Multiplying an extra $2\times$
for INT8 into a cycle product is a double-count. Dropping it lands the mechanical
ladder at $35$--$37\times$, within $3$--$12\%$ of the headline.
\item \textbf{The software rows use a $\sim1.3\times$ smaller naive baseline}
(the scale-$T_{\max}$ naive) than the config-0 naive against which the headline
is measured, so the two families of factors are not on the same numerator.
\end{enumerate}
\Cref{eq:twofactor} is the form that reconciles \emph{exactly} because $f_1$ and
$f_2$ share a single naive numerator and a single cycle basis. As a naming
caution: the same fixed-silicon scheduling recovery is $1.48$--$1.78\times$
against the config-0 naive (used above) and $2.14$--$2.51\times$ against the
grid's scale-$T_{\max}$ naive (the fixed-hardware grid convention); the two
differ only by the naive-baseline
definition, and must never be crossed within one ratio.

\paragraph{The provisioning factor is sub-linear (the anti-``just add MACs''
result).} $f_2$ is not free scaling. Measured as delivered throughput per
provisioned MAC in an equal-area / equal-TDP resource-matched ablation,
$20.0$--$27.0\times$ more MACs realize only $\approx16\times$ provisioning
throughput, a linearity of $\mathbf{0.59}$--$\mathbf{0.81}$ (Llama-3 70B is
worst, $27\times$ MACs $\to$ $16.0\times$). Equivalently, the fixed-family MAC
leg from $65{,}536$ to $1{,}048{,}576$ MACs (a $16\times$ increase) yields only
$14.0$--$15.7\times$ cycle reduction, because a residual memory floor remains.
Added silicon buys sub-linearly, which is exactly why capturing the free
scheduling factor $f_1$ first (and provisioning MACs only up to the point of
diminishing return) is the efficient order of operations.
\fi

% ---------------------------------------------------------------------------
\subsection{One-Decision-at-a-Time Sensitivity Sweeps}
\label{app:deficit:sweeps}

We now hold the design fixed and sweep a single decision, reporting the objective
(LP cycles) sensitivity. Two decisions have dedicated measured sweeps: global buffer and MAC count;
Table~\ref{tab:deficit_sweepsummary} summarizes them.

\paragraph{(i) Global-buffer size $C_{\mathrm{GM}}$: measured, near-inert in
the compute-bound regime.} \Cref{tab:deficit_cgm} sweeps $C_{\mathrm{GM}}$ from
$1$ to $64$\,MiB for DeepSeek-V3 on TPU-v3 ($900$\,GB/s), for the joint (free-$y$)
LP and for a fixed-tile FAST-subset LP. The joint LP is \emph{flat}: its objective
moves $0.39\%$ across a $64\times$ buffer sweep ($1.0723\times10^{12}\to
1.0682\times10^{12}$ cycles), and is bit-identical from $8$\,MiB upward. By
contrast the fixed-tile schedule \emph{needs} the buffer, improving $1.18\times$
over the same sweep. The joint formulation substitutes tile choice for buffer
capacity: when tiling is free, a larger global buffer buys almost nothing. The
tile-dependent capacity extension is likewise inert here: it matches the retained joint model
to within $0.08\%$ across the full $C_{\mathrm{GM}}\in\{4,\dots,64\}$\,MiB sweep.
Under bandwidth pressure ($64$--$128$\,GB/s) the buffer regains a modest decision
($\approx1.19\times$ for Llama-1 13B over $8\to64$\,MiB), but \emph{only through
placement}: with placement frozen (row A) the objective is constant across
$C_{\mathrm{GM}}$, so $C_{\mathrm{GM}}$ helps exclusively by enabling
time-indexed residency, never on its own. The blueprints adopt
$C_{\mathrm{GM}}=64$\,MiB.

\begin{table}[ht!]
\centering
\caption{Global-buffer sweep, DeepSeek-V3 on TPU-v3 ($900$\,GB/s, native cache
convention, HiGHS-Optimal). The joint (free-$y$) LP is essentially
$C_{\mathrm{GM}}$-insensitive ($0.39\%$ over a $64\times$ sweep); the fixed-tile
FAST-subset LP is not ($1.18\times$). Objective in LP cycles.}
\label{tab:deficit_cgm}
\small
\resizebox{\linewidth}{!}{%
\begin{tabular}{@{}lcccccccc@{}}
\toprule
$C_{\mathrm{GM}}$ (MiB) & 1 & 2 & 4 & 8 & 16 & 32 & 64 & sweep ($1\!\to\!64$) \\
\midrule
Joint (free-$y$) & 1.0723 & 1.0699 & 1.0690 & 1.0682 & 1.0682
& 1.0682 & 1.0682 & $\times10^{12}$; $1.004\times$ \\
Fixed-tile (FAST-subset)   & 1.4627 & 1.4459 & 1.4327 & 1.4161 & 1.3863
& 1.3364 & 1.2371 & $\times10^{12}$; $1.18\times$ \\
\bottomrule
\end{tabular}%
}
\end{table}

\paragraph{(ii) Peak throughput / MAC-unit count.} The total-MAC provisioning sweep is measured: raising the PE
grid from $65{,}536$ to $1{,}048{,}576$ MACs (a $16\times$ increase) reduces the
objective by $14.0$--$15.7\times$ (Llama-1 13B $15.72\times$, DeepSeek-V3
$14.98\times$, Llama-3 70B $14.04\times$, Llama-1 65B $14.02\times$), and the
provisioning factor to the full co-designed chip is $f_2=20.0$--$26.9\times$
both sub-linear.
This is the single largest decision on the objective.

\paragraph{(iii) DRAM bandwidth: measured, inert in the compute-bound regime.}
Sweeping bandwidth from $900$ to $3{,}350$\,GB/s at the fixed $65{,}536$-MAC
point moves the objective by $\le0.25\%$ (BW leg $=1.000$--$1.0025\times$ for
all four LLMs): at the deployed prefill-like batch these workloads are
compute-bound, so bandwidth beyond the floor buys nothing. Bandwidth becomes the
dominant decision only in the bandwidth-starved regime: at $64$\,GB/s (the compact
NPU) the joint-vs-naive fixed-hardware gain is $5.9$--$31.4\times$, because there
the roofline is memory-bound and residency / placement recover a much larger
share. Thus bandwidth sensitivity is regime-dependent: near-zero when
compute-bound, first-order when memory-bound.

\paragraph{(iv) Interconnect: on-chip NoC and inter-device topology.}
\emph{On-chip}, the NoC / crossbar radix and link count set the PE-to-SRAM
routing capacity; this is the resource bounded by the crossbar / NoC-routing
gate of \cref{app:gates} (gate 4, $N_{\mathrm{NoC}}=n_{\mathrm{PEs}}\times
n_{\mathrm{banks}}\le5\times10^{6}$ cells), and its provisioning follows the
network-on-package accounting of Simba~\citep{shao2019simba}. In the single-chip
compute-bound regime the on-chip NoC is provisioned to the compute floor and is
not the binding resource.
\emph{Across devices}, in the multi-GPU / mixture-of-experts regime the
inter-device topology (radix and link count of the NVLink-class fabric) binds
only when experts are dispatched \emph{across} devices; when the active-expert
set is package-resident, as in every blueprint here, the communication term is
zero.

\begin{table}[ht!]
\centering
\caption{Single-decision sensitivity summary (compute-bound LLM regime, TPU-v3
operating point). ``Objective sensitivity'' is the LP-cycle response to the
stated sweep.}
\label{tab:deficit_sweepsummary}
\small
\begin{tabularx}{\linewidth}{@{}p{.22\linewidth}p{.28\linewidth}X@{}}
\toprule
\textbf{Decision} & \textbf{Sweep} & \textbf{Objective sensitivity} \\
\midrule
Global buffer $C_{\mathrm{GM}}$ & $1\to64$\,MiB (BW$=900$)
& $1.004\times$ (joint); $1.18\times$ (fixed-tile) \\
& under BW pressure ($\le128$) & $\approx1.19\times$, via placement only \\
MAC count (total) & $65{,}536\to1{,}048{,}576$ ($16\times$)
& $14.0$--$15.7\times$ (sub-linear); $f_2=20.0$--$26.9\times$ \\
DRAM bandwidth & $900\to3{,}350$\,GB/s & $\le1.003\times$ (compute-bound) \\
& $64$\,GB/s (NPU) & $5.9$--$31.4\times$ (memory-bound) \\
\bottomrule
\end{tabularx}
\end{table}

% ---------------------------------------------------------------------------
